\documentclass[11pt]{article}
\usepackage[margin=1in]{geometry}
\usepackage{amsmath,amssymb,amsthm,amscd,mathtools,array,booktabs}
\usepackage[T1]{fontenc}
\usepackage{lmodern}
\usepackage{microtype}
\usepackage[hidelinks]{hyperref}
\usepackage[nameinlink,capitalise]{cleveref}

\newtheorem{theorem}{Theorem}[section]
\newtheorem{proposition}[theorem]{Proposition}
\newtheorem{lemma}[theorem]{Lemma}
\newtheorem{corollary}[theorem]{Corollary}
\theoremstyle{definition}
\newtheorem{definition}[theorem]{Definition}
\theoremstyle{remark}
\newtheorem{remark}[theorem]{Remark}

\newcommand{\C}{\mathbb C}
\newcommand{\Hh}{\mathbb H}
\newcommand{\R}{\mathbb R}
\newcommand{\Z}{\mathbb Z}
\newcommand{\CP}{\mathbb {CP}}
\newcommand{\Herm}{\operatorname{Herm}}
\newcommand{\Mat}{\operatorname{Mat}}
\newcommand{\rank}{\operatorname{rank}}
\newcommand{\tr}{\operatorname{tr}}
\newcommand{\Ran}{\operatorname{Ran}}
\newcommand{\adj}{\operatorname{adj}}
\newcommand{\diag}{\operatorname{diag}}
\newcommand{\Sp}{\operatorname{Sp}}
\newcommand{\id}{\operatorname{id}}
\newcommand{\cC}{\mathcal C}
\newcommand{\cJ}{\mathcal J}
\newcommand{\cN}{\mathcal N}
\newcommand{\cO}{\mathcal O}
\newcommand{\cS}{\mathcal S}
\newcommand{\cK}{\mathcal K}
\newcommand{\cG}{\mathcal G}
\newcommand{\eps}{\varepsilon}
\newcommand{\bj}{\mathbf j}
\newcommand{\bk}{\mathbf k}
\newcommand{\e}{\mathrm e}

\title{The Exact \texorpdfstring{$4+4+4+1$}{4+4+4+1} Continuation:\\
A Thirteen-Dimensional Quaternionic Peirce Slice, Its Dirac--Lorentz Cubic, and Its Rank-Defect \texorpdfstring{$S^{12}$}{S12}}
\author{Working theorem/proof module}
\date{4 September 2026}

\begin{document}
\maketitle

\begin{abstract}
The previously established marked observation matrix
\[
 P=\begin{pmatrix}1&-1&0\\0&1&-1\end{pmatrix}
\]
splits $\Herm_3(\C)$ as a four-dimensional Lorentz vector, one four-real-dimensional Weyl spinor, and a scalar.  This note extends the same map over the quaternions without replacing its basis or its coefficients.  Fixing the observed block to remain in $\Herm_2(\C)$ gives the pullback
\[
 \cJ_{13,P}:=(\cO_P^{\Hh})^{-1}(\Herm_2(\C))\subset\Herm_3(\Hh),
\]
a real thirteen-dimensional space with the exact decomposition
\[
 \cJ_{13,P}\cong\Herm_2(\C)\oplus\C^2\oplus\C^2\oplus\R.
\]
Its restricted Moore determinant is the cubic
\[
 \cN_{13}(G,z,w,\lambda)
 =\lambda\det G-z^*\adj(G)z-w^*\adj(G)w.
\]
After charge conjugating the second spinor this becomes the Lorentz-invariant vector--Dirac-spinor--scalar expression
\[
 \lambda\det G-\psi_L^*\adj(G)\psi_L-\psi_R^*G\psi_R.
\]
We classify every positive-semidefinite fibre, prove invariance under
$(SL_2(\C)\times SU(2))/\{\pm(1,1)\}$, identify an $\CP^1$-family of determinant-preserving nine-dimensional slices, and show that the old compactified rank defect is the fixed unknotted $S^8\subset S^{12}$.  The complement is $S^3\times\R^9$, so the added four-dimensional spinor has the unit-quaternion linking sphere predicted by the coordinate count.  The slice is not a Jordan subalgebra; an explicit product leaves it, and its Jordan closure is all of $\Herm_3(\Hh)$.  We also compare it, on the identical coordinate space, with the other natural thirteen-dimensional complex $4\times4$ determinant model and prove the exact correction
\[
 \det\begin{pmatrix}G&Z\\Z^*&\lambda I_2\end{pmatrix}
 =\lambda\cN_{13}(G,Z,\lambda)+|\det Z|^2.
\]
Thus the quaternionic model is distinguished by preserving the old cubic exactly, while the complex rank-four model carries an additional Pluecker term.  No toroidal filling or quantum field theory is identified by the dimension count alone.
\end{abstract}

\tableofcontents

\section{The retained observation map and the quaternionic pullback}

Write the quaternions as
\begin{equation}\label{eq:H-decomp}
 \Hh=\C\oplus\C\bj,
 \qquad \bj^2=-1,
 \qquad \bj a=\overline a\,\bj\quad(a\in\C).
\end{equation}
For $q=a+b\bj$, quaternionic conjugation is
$\overline q=\overline a-b\bj$.

Retain the exact marked data
\begin{equation}\label{eq:P-data}
 P=\begin{pmatrix}1&-1&0\\0&1&-1\end{pmatrix},
 \qquad
 u=\begin{pmatrix}1\\1\\1\end{pmatrix},
 \qquad
 C=PP^*=\begin{pmatrix}2&-1\\-1&2\end{pmatrix},
\end{equation}
\begin{equation}\label{eq:R-data}
 R_P=P^*C^{-1}
 =\frac13\begin{pmatrix}2&1\\-1&1\\-1&-2\end{pmatrix},
 \qquad
 \Pi_u=\frac{uu^*}{3},
 \qquad X=(R_P\;u).
\end{equation}
Then
\begin{equation}\label{eq:split}
 PR_P=I_2,
 \quad R_PP=I_3-\Pi_u,
 \quad u^*R_P=0,
 \quad Pu=0,
 \quad \det X=1,
 \quad X^{-1}=\begin{pmatrix}P\\u^*/3\end{pmatrix}.
\end{equation}

Extend the observation map, using the same real matrix $P$, to
\begin{equation}\label{eq:OH}
 \cO_P^{\Hh}:\Herm_3(\Hh)\longrightarrow\Herm_2(\Hh),
 \qquad \cO_P^{\Hh}(H)=PHP^*.
\end{equation}
Fix the displayed complex subalgebra $\C\subset\Hh$ and define
\begin{equation}\label{eq:J13-pullback}
 \boxed{
 \cJ_{13,P}:=(\cO_P^{\Hh})^{-1}(\Herm_2(\C)).}
\end{equation}
Thus $\cJ_{13,P}$ is defined by a pullback square, not by dimension alone:
\begin{equation}\label{eq:pullback-square}
\begin{CD}
 \cJ_{13,P}@>>>\Herm_3(\Hh)\\
 @V{\cO_P}VV @VV{\cO_P^{\Hh}}V\\
 \Herm_2(\C)@>>>\Herm_2(\Hh).
\end{CD}
\end{equation}

\begin{theorem}[Unique $4+4+4+1$ coordinates]\label{thm:coordinates}
Every $H\in\cJ_{13,P}$ has one and only one expression
\begin{equation}\label{eq:marked-block}
 H=X
 \begin{pmatrix}
  G&\psi\\ \psi^*&\lambda
 \end{pmatrix}X^*,
 \qquad
 G\in\Herm_2(\C),\quad
 \psi\in\Hh^2,\quad
 \lambda\in\R.
\end{equation}
Writing
\begin{equation}\label{eq:psi-zw}
 \psi=z+w\bj,
 \qquad z,w\in\C^2,
\end{equation}
the extraction formulae are
\begin{equation}\label{eq:extraction}
 G=PHP^*,
 \qquad
 z+w\bj=\frac13PHu,
 \qquad
 \lambda=\frac19u^*Hu.
\end{equation}
Consequently
\begin{equation}\label{eq:4441}
 \boxed{
 \cJ_{13,P}
 \cong_{\R}
 \underbrace{\Herm_2(\C)}_{4}
 \oplus
 \underbrace{\C^2}_{4}
 \oplus
 \underbrace{\C^2}_{4}
 \oplus
 \underbrace{\R}_{1}.}
\end{equation}
The old space $\Herm_3(\C)$ is exactly the subspace $w=0$.
\end{theorem}

\begin{proof}
For an arbitrary quaternionic Hermitian matrix, conjugation by $X^{-1}$ gives
\[
 X^{-1}HX^{-*}
 =\begin{pmatrix}
 PHP^*&PHu/3\\u^*HP^*/3&u^*Hu/9
 \end{pmatrix}.
\]
Membership in $\cJ_{13,P}$ says precisely that the upper-left block is complex Hermitian.  The remaining off-diagonal block is an arbitrary vector in $\Hh^2$, and \eqref{eq:H-decomp} gives its unique pair $(z,w)$.  The dimension is $4+8+1=13$.  If $w=0$, every block is complex, and because $X$ is real this is exactly $\Herm_3(\C)$; conversely every complex Hermitian matrix has $w=0$.
\end{proof}

\section{The restricted Moore determinant}

Use the complex $2\times2$ representation
\begin{equation}\label{eq:rhoH}
 \rho(a+b\bj)=
 \begin{pmatrix}a&b\\-\overline b&\overline a\end{pmatrix}.
\end{equation}
For a quaternionic Hermitian matrix $M$, its Moore determinant $\det_{\Hh}M$ is the real polynomial normalized by $\det_{\Hh}I=1$ and
\begin{equation}\label{eq:Moore-square}
 \det_{\C}\rho(M)=(\det_{\Hh}M)^2.
\end{equation}

\begin{lemma}[Quaternionic form splitting]\label{lem:qsplit}
If $A\in\Herm_2(\C)$ and $\psi=z+w\bj\in\Hh^2$, then
\begin{equation}\label{eq:qsplit}
 \boxed{\psi^*A\psi=z^*Az+w^*Aw\in\R.}
\end{equation}
\end{lemma}

\begin{proof}
Expand with $\bj a=\overline a\bj$.  The two complex scalar parts are $z^*Az$ and $w^*Aw$.  The coefficient of $\bj$ is
$z^*Aw-w^{\mathsf T}\overline A\,\overline z$, which vanishes because $A=A^*$.  Both scalar terms are real.
\end{proof}

\begin{theorem}[The thirteen-dimensional cubic]\label{thm:cubic}
For $H$ with coordinates $(G,z,w,\lambda)$, the restricted Moore determinant is
\begin{equation}\label{eq:N13}
 \boxed{
 \cN_{13}(G,z,w,\lambda)
 :=\det_{\Hh}H
 =\lambda\det_{\C}G
  -z^*\adj(G)z
  -w^*\adj(G)w.}
\end{equation}
Equivalently, with $Z=(z\;w)\in\Mat_{2\times2}(\C)$,
\begin{equation}\label{eq:N13-trace}
 \cN_{13}(G,Z,\lambda)
 =\lambda\det G-\tr\bigl(Z^*\adj(G)Z\bigr).
\end{equation}
On $w=0$ this is exactly the old nine-dimensional cubic
\begin{equation}\label{eq:N9}
 \cN_9(G,z,\lambda)=\lambda\det G-z^*\adj(G)z,
\end{equation}
with no extra factor.
\end{theorem}

\begin{proof}
Since $\det X=1$, quaternionic congruence by $X$ leaves the Moore determinant unchanged.  For invertible $G$, the quaternionic Schur formula gives
\[
 \det_{\Hh}
 \begin{pmatrix}G&\psi\\\psi^*&\lambda\end{pmatrix}
 =\det G\,(\lambda-\psi^*G^{-1}\psi)
 =\lambda\det G-\psi^*\adj(G)\psi.
\]
Both sides are polynomial, so the identity holds for singular $G$ as well.  Apply \cref{lem:qsplit} with $A=\adj(G)$.
\end{proof}

\subsection{Dirac form of the same cubic}

Put
\begin{equation}\label{eq:epsilon}
 \epsilon_2=\begin{pmatrix}0&1\\-1&0\end{pmatrix}
\end{equation}
and write the second left-handed spinor as the charge conjugate of a right-handed one:
\begin{equation}\label{eq:charge-conj}
 w=\epsilon_2\overline{\psi_R},
 \qquad z=\psi_L.
\end{equation}
For every $G\in\Herm_2(\C)$,
\begin{equation}\label{eq:epsilon-adj}
 \epsilon_2^*\adj(G)\epsilon_2=G^{\mathsf T}.
\end{equation}

\begin{corollary}[Vector--Dirac-spinor--scalar cubic]\label{cor:Dirac}
In the variables $(G,\psi_L,\psi_R,\lambda)$,
\begin{equation}\label{eq:Dirac-cubic}
 \boxed{
 \cN_{13}
 =\lambda\det G
 -\psi_L^*\adj(G)\psi_L
 -\psi_R^*G\psi_R.}
\end{equation}
Under $g\in SL_2(\C)$,
\begin{equation}\label{eq:chiral-action}
 G\mapsto gGg^*,
 \qquad
 \psi_L\mapsto g\psi_L,
 \qquad
 \psi_R\mapsto(g^*)^{-1}\psi_R,
 \qquad
 \lambda\mapsto\lambda.
\end{equation}
Thus the two four-real-dimensional summands may be read either as two left Weyl spinors or, after the exact map \eqref{eq:charge-conj}, as one Dirac spinor.
\end{corollary}

\begin{proof}
Equation \eqref{eq:epsilon-adj} gives
\[
 w^*\adj(G)w
 =\psi_R^{\mathsf T}G^{\mathsf T}\overline{\psi_R}
 =\psi_R^*G\psi_R.
\]
The identity $g\epsilon_2=\epsilon_2(g^{\mathsf T})^{-1}$ proves that $w\mapsto gw$ is equivalent to the right-handed transformation in \eqref{eq:chiral-action}.
\end{proof}

\section{The complete positive cone and all observation fibres}

Let
\begin{equation}\label{eq:C13}
 \cC_{13,P}=\cJ_{13,P}\cap\Herm_3(\Hh)_{\succeq0}.
\end{equation}

\begin{theorem}[Exact positivity and rank]\label{thm:positivity}
A point $(G,z,w,\lambda)$ lies in $\cC_{13,P}$ if and only if
\begin{equation}\label{eq:positivity}
 \boxed{
 G\succeq0,
 \qquad z,w\in\Ran G,
 \qquad
 \lambda\ge z^*G^\dagger z+w^*G^\dagger w.}
\end{equation}
Put
\begin{equation}\label{eq:s-schur}
 s=\lambda-z^*G^\dagger z-w^*G^\dagger w\ge0.
\end{equation}
Then the quaternionic rank is
\begin{equation}\label{eq:rank}
 \boxed{
 \rank_{\Hh}H
 =\rank_{\C}G+
 \begin{cases}1,&s>0,\\0,&s=0.\end{cases}}
\end{equation}
For fixed $G\succeq0$ of rank $r$, the positive fibre of $\cO_P$ is
\begin{equation}\label{eq:fibre}
 \left\{
 (z,w,s):z,w\in\Ran G,\ s\ge0
 \right\}
 \cong\C^{2r}\times[0,\infty),
\end{equation}
of real dimension $4r+1$.  In particular,
\begin{equation}\label{eq:zero-fibre}
 \boxed{
 \cO_P^{-1}(0)\cap\cC_{13,P}
 =\{\lambda uu^*: \lambda\ge0\}.}
\end{equation}
\end{theorem}

\begin{proof}
Work in the block basis $X$.  A positive block matrix must have $G\succeq0$.  Positivity on vectors from $\ker G\oplus\Hh$ forces $\psi=z+w\bj$ to lie in the quaternionic range of $G$.  Since $G$ is complex, this is equivalent to $z,w\in\Ran_{\C}G$.  Put $v=G^\dagger\psi$.  The exact quaternionic congruence
\[
 \begin{pmatrix}I_2&0\\-v^*&1\end{pmatrix}
 \begin{pmatrix}G&\psi\\\psi^*&\lambda\end{pmatrix}
 \begin{pmatrix}I_2&-v\\0&1\end{pmatrix}
 =\begin{pmatrix}G&0\\0&\lambda-\psi^*G^\dagger\psi\end{pmatrix}
\]
proves positivity and rank.  By \cref{lem:qsplit}, the last scalar is exactly \eqref{eq:s-schur}.  If $G=0$, then $z=w=0$, leaving the ray \eqref{eq:zero-fibre} after returning through $X$.
\end{proof}

\begin{corollary}[Boundary strata]\label{cor:strata}
The open positive-definite stratum has real dimension $13$.  The generic rank-two boundary stratum is given by $G>0$ and $s=0$ and has dimension $12$.  The rank-one locus has top dimension $7$, obtained from $\rank G=1$ and $s=0$; the ray \eqref{eq:zero-fibre} is contained in its closure.
\end{corollary}

\section{Lorentz and internal symmetry}

For $g\in SL_2(\C)$ retain
\begin{equation}\label{eq:TP}
 T_P(g)=X(g\oplus1)X^{-1}=R_PgP+\Pi_u.
\end{equation}
For $k\in SU(2)$ act on the two-spinor matrix $Z=(z\;w)$ on the right.

\begin{theorem}[Exact symmetry group]\label{thm:symmetry}
The action
\begin{equation}\label{eq:group-action}
 (g,k)\cdot(G,Z,\lambda)
 =\bigl(gGg^*,\ gZk,\ \lambda\bigr)
\end{equation}
preserves $\cC_{13,P}$ and $\cN_{13}$.  Its kernel is
\begin{equation}\label{eq:group-kernel}
 \{(I_2,I_2),(-I_2,-I_2)\},
\end{equation}
so the effective group is
\begin{equation}\label{eq:effective-group}
 \boxed{(SL_2(\C)\times SU(2))/\mu_2.}
\end{equation}
The left factor is exactly the previous marked Lorentz action.  The right factor is right multiplication of $\psi\in\Hh^2$ by unit quaternions and rotates the two complex Weyl summands.
\end{theorem}

\begin{proof}
For $g\in SL_2(\C)$,
\[
 \adj(gGg^*)=(g^*)^{-1}\adj(G)g^{-1},
\]
so both spinor quadratic terms and $\det G$ are invariant.  For $k\in SU(2)$,
\[
 \tr\bigl((Zk)^*\adj(G)(Zk)\bigr)
 =\tr\bigl(k^*Z^*\adj(G)Zk\bigr)
 =\tr\bigl(Z^*\adj(G)Z\bigr).
\]
The positivity criterion \eqref{eq:positivity} has the same invariance.  The only elements acting trivially on every $G$ are $g=\pm I_2$, and triviality on every $Z$ then forces $k=g$.  The quaternionic realization follows from the usual identification of $SU(2)$ with unit quaternions and the decomposition \eqref{eq:H-decomp}.
\end{proof}

\section{The old nine-dimensional object inside the new one}

Define
\begin{align}
 \iota_9(G,z,\lambda)&=(G,z,0,\lambda),\label{eq:iota9}\\
 r_9(G,z,w,\lambda)&=(G,z,\lambda).\label{eq:r9}
\end{align}

\begin{theorem}[Split positive continuation]\label{thm:split-continuation}
The map $\iota_9$ is an injective, determinant-preserving, positivity-preserving linear map
\begin{equation}\label{eq:iota-cones}
 \iota_9:\Herm_3(\C)_{\succeq0}\hookrightarrow\cC_{13,P}.
\end{equation}
The coordinate retraction $r_9$ maps $\cC_{13,P}$ into $\Herm_3(\C)_{\succeq0}$ and satisfies
\begin{equation}\label{eq:split-r}
 r_9\iota_9=\id.
\end{equation}
Their determinants are related by the exact correction
\begin{equation}\label{eq:det-correction}
 \boxed{
 \cN_9(r_9H)-\cN_{13}(H)
 =w^*\adj(G)w\ge0
 \quad(H\in\cC_{13,P}).}
\end{equation}
Both maps are proper on their closed positive cones and therefore extend to the one-point compactifications
\begin{equation}\label{eq:ball-split}
 B^9\xrightarrow{\ \iota_9^+\ }B^{13}
 \xrightarrow{\ r_9^+\ }B^9,
 \qquad r_9^+\iota_9^+=\id.
\end{equation}
The retraction does not restrict to a map from the new rank-defect boundary to the old one: if $G=I_2$, $z=0$, $w\ne0$, and $\lambda=\|w\|^2$, then $\cN_{13}=0$ but $\cN_9=\lambda>0$.
\end{theorem}

\begin{proof}
The first statements follow immediately from \cref{thm:cubic,thm:positivity}.  Since $G\succeq0$ implies $\adj(G)\succeq0$, the correction is nonnegative.  To prove properness of $r_9$ on the cone, let $(G,z,\lambda)$ range over a compact set.  If $g_i$ are the eigenvalues of $G$, \eqref{eq:positivity} gives, in a $G$-eigenbasis,
\[
 |w_i|^2\le \lambda g_i.
\]
Thus $w$ is bounded; the fibre is also closed.  Properness of $\iota_9$ is immediate.  The displayed counterexample proves the final assertion and supplies the exact map relating the two boundaries when a boundary-to-boundary retraction fails.
\end{proof}

\subsection{A family of old slices}

For a line $L=[\ell]\in\CP^1$, choose $\ell\in\C^2$ with $\ell^*\ell=1$ and set
\begin{equation}\label{eq:J9L}
 \cJ_9(L)=\{(G,Z,\lambda):Z=\psi\ell^*,\ \psi\in\C^2\}.
\end{equation}
This subspace is independent of the phase of $\ell$.

\begin{proposition}[$\CP^1$-family of determinant-preserving slices]\label{prop:CP1-family}
Each $\cJ_9(L)$ is a nine-dimensional copy of $\Herm_3(\C)$, and
\begin{equation}\label{eq:N-on-L}
 \cN_{13}(G,\psi\ell^*,\lambda)
 =\lambda\det G-\psi^*\adj(G)\psi.
\end{equation}
The internal $SU(2)$ acts transitively on this family, with stabilizer $U(1)$, so its parameter space is $SU(2)/U(1)=\CP^1$.  The locus contained in at least one such slice is exactly $\rank_{\C}Z\le1$.
\end{proposition}

\begin{proof}
One has
\[
 \tr\bigl((\psi\ell^*)^*\adj(G)(\psi\ell^*)\bigr)
 =\tr\bigl(\ell\psi^*\adj(G)\psi\ell^*\bigr)
 =\psi^*\adj(G)\psi.
\]
The remaining statements are the standard transitive action on complex lines and the rank-one factorization of a $2\times2$ complex matrix.
\end{proof}

\section{The rank-defect sphere is exactly \texorpdfstring{$S^{12}$}{S12}}

Let
\begin{equation}\label{eq:boundary13}
 \partial\cC_{13,P}
 =\{H\in\cC_{13,P}:\det_{\Hh}H=0\}
\end{equation}
and let
\begin{equation}\label{eq:tracefree13}
 \cJ_{13,P}^0=\{K\in\cJ_{13,P}:\tr K=0\}.
\end{equation}
This is a real vector space of dimension $12$.

\begin{theorem}[Spectral-floor coordinates]\label{thm:spectral-floor}
For quaternionic Hermitian $K$, let $\lambda_{\min}(K)$ denote its least real eigenvalue.  The maps
\begin{align}
 \Phi_{13}(K)&=K-\lambda_{\min}(K)I_3,
 &&K\in\cJ_{13,P}^0,\label{eq:Phi13}\\
 \Psi_{13}(H)&=H-\frac{\tr H}{3}I_3,
 &&H\in\partial\cC_{13,P},\label{eq:Psi13}
\end{align}
are inverse positively homogeneous homeomorphisms.  More generally,
\begin{equation}\label{eq:halfspace-homeo}
 \cC_{13,P}\xrightarrow{\cong}
 \cJ_{13,P}^0\times[0,\infty),
 \qquad
 H\longmapsto
 \left(H-\frac{\tr H}{3}I_3,
       \frac{\tr H}{3}+\lambda_{\min}\left(H-\frac{\tr H}{3}I_3\right)
 \right).
\end{equation}
Consequently
\begin{equation}\label{eq:compact-topology}
 \boxed{
 \cC_{13,P}^+\cong B^{13},
 \qquad
 (\partial\cC_{13,P})^+\cong S^{12}.}
\end{equation}
\end{theorem}

\begin{proof}
The identity matrix lies in $\cJ_{13,P}$, so both shifts remain in the slice.  If $K$ is trace-free, $K-\lambda_{\min}(K)I_3$ is positive semidefinite and singular.  If $H$ is positive semidefinite and singular, then
\[
 \lambda_{\min}\left(H-\frac{\tr H}{3}I_3\right)
 =-\frac{\tr H}{3},
\]
so the maps are inverse.  Continuity follows from continuity of ordered Hermitian eigenvalues.  Writing $H=K+rI_3$ shows that positivity is equivalent to $r\ge-\lambda_{\min}(K)$; the second coordinate in \eqref{eq:halfspace-homeo} is exactly the nonnegative excess.  The one-point compactifications of $\R^{12}\times[0,\infty)$ and its boundary $\R^{12}$ are $B^{13}$ and $S^{12}$.
\end{proof}

\begin{theorem}[The old $S^8$ and its quaternionic link]\label{thm:S8-in-S12}
The involution
\begin{equation}\label{eq:involution}
 \sigma(G,z,w,\lambda)=(G,z,-w,\lambda)
\end{equation}
preserves $\cN_{13}$ and $\cC_{13,P}$.  Its fixed space is $\Herm_3(\C)$, and its fixed set on the compactified rank defect is the previously constructed $S^8$.  Under the spectral-floor coordinates,
\begin{equation}\label{eq:pair}
 \boxed{
 \bigl((\partial\cC_{13,P})^+,
       (\partial\Herm_3(\C)_{\succeq0})^+\bigr)
 \cong(S^{12},S^8)}
\end{equation}
as the standard unknotted linear pair.  Its complement is
\begin{equation}\label{eq:complement}
 \boxed{S^{12}\setminus S^8\cong S^3\times\R^9.}
\end{equation}
The $S^3$ is the unit sphere in the added $w\in\C^2\cong\R^4$ direction and is the unit-quaternion group $\Sp(1)\cong SU(2)$.
\end{theorem}

\begin{proof}
The determinant and positivity formulae are even in $w$.  The fixed locus $w=0$ is $\Herm_3(\C)$ by \cref{thm:coordinates}.  The trace-free pair is the linear pair
\[
 (\cJ_{13,P}^0,\Herm_3^0(\C))\cong(\R^{12},\R^8).
\]
One-point compactification gives the standard pair.  Removing the fixed set gives
\[
 \R^8\times(\R^4\setminus\{0\})
 \cong\R^8\times S^3\times\R
 \cong S^3\times\R^9.
\]
\end{proof}

\section{It is a generating cubic slice, not a Jordan subalgebra}

The ambient quaternionic Euclidean Jordan algebra has the Peirce decomposition
\begin{equation}\label{eq:H3H-Peirce}
 \Herm_3(\Hh)
 \cong
 \underbrace{\Herm_2(\Hh)}_{6}
 \oplus
 \underbrace{\Hh^2}_{8}
 \oplus
 \underbrace{\R}_{1}.
\end{equation}
The inclusion $\Herm_2(\C)\subset\Herm_2(\Hh)$ selects the $(1,3)$ Lorentz subspace inside the full $(1,5)$ spin factor and gives exactly $\cJ_{13,P}$.

\begin{theorem}[Exact Jordan completion]\label{thm:Jordan-closure}
The space $\cJ_{13,P}$ is not closed under the Jordan product
$A\circ B=(AB+BA)/2$.  In the block basis, take
\[
 \psi=\binom10,
 \qquad
 \phi=\binom0{\bj},
 \qquad
 A=\begin{pmatrix}0&\psi\\\psi^*&0\end{pmatrix},
 \qquad
 B=\begin{pmatrix}0&\phi\\\phi^*&0\end{pmatrix}.
\]
Then the upper-left block of $A\circ B$ is
\begin{equation}\label{eq:nonclosure}
 \frac12
 \begin{pmatrix}0&-\bj\\\bj&0\end{pmatrix},
\end{equation}
which is not in $\Herm_2(\C)$.  Replacing $\bj$ by $\bk$ produces the other missing real direction.  Therefore
\begin{equation}\label{eq:Jordan-closure}
 \boxed{
 \operatorname{Jord}\langle\cJ_{13,P}\rangle
 =\Herm_3(\Hh).}
\end{equation}
The exact completion is
\begin{equation}\label{eq:15-completion}
 4+4+4+1
 \quad\longrightarrow\quad
 6+4+4+1=15,
\end{equation}
where the two added coordinates complete the Lorentz vector from $(1,3)$ to $(1,5)$; the two spinors and scalar are unchanged.
\end{theorem}

\begin{proof}
The displayed product is a direct multiplication.  The slice already contains the full $\Hh^2$ Peirce half-space, the scalar, and the complex four-dimensional part of $\Herm_2(\Hh)$.  The two products with $\bj$ and $\bk$ generate the two missing off-diagonal components of $\Herm_2(\Hh)$.  Hence the Jordan algebra generated is the full ambient algebra.
\end{proof}

\section{The other thirteen-dimensional determinant and the exact comparison}

The same coordinate vector space has a second natural realization as a complex Hermitian $4\times4$ slice:
\begin{equation}\label{eq:K4}
 \cK_{13}(G,Z,\lambda)
 =\begin{pmatrix}G&Z\\Z^*&\lambda I_2\end{pmatrix}
 \in\Herm_4(\C).
\end{equation}
It has the same $4+8+1$ coordinate count and the same
$SL_2(\C)\times SU(2)$ action, but its determinant is quartic.

\begin{theorem}[Quartic--cubic comparison]\label{thm:quartic-comparison}
For all $(G,Z,\lambda)$,
\begin{align}
 \det_{\C}\cK_{13}(G,Z,\lambda)
 &=\lambda^2\det G
 -\lambda\tr\bigl(Z^*\adj(G)Z\bigr)
 +|\det Z|^2\label{eq:D4}\\
 &=\boxed{\lambda\cN_{13}(G,Z,\lambda)+|\det Z|^2.}\label{eq:D4N13}
\end{align}
On the old slice $Z=(z\;0)$,
\begin{equation}\label{eq:quartic-old}
 \det\cK_{13}=\lambda\cN_9,
 \qquad
 \cN_{13}=\cN_9.
\end{equation}
Thus the quaternionic slice is the one that preserves the preceding cubic exactly; the complex rank-four slice adds both a factor $\lambda$ and the Pluecker correction $|\det Z|^2$.

The coordinate identity maps the quaternionic positive cone into the complex one:
\begin{equation}\label{eq:cone-inclusion}
 \cC_{13,P}\hookrightarrow
 \{(G,Z,\lambda):\cK_{13}(G,Z,\lambda)\succeq0\}.
\end{equation}
It is not onto.  For example,
\begin{equation}\label{eq:not-onto-example}
 G=I_2,
 \qquad Z=I_2,
 \qquad \lambda=1
\end{equation}
gives a positive-semidefinite complex $4\times4$ block, while the quaternionic condition would require $1\ge\tr(I_2)=2$.
\end{theorem}

\begin{proof}
For invertible $G$ or nonzero $\lambda$, the block determinant formula gives \eqref{eq:D4}; polynomial continuation gives it everywhere.  The last term is
$\det(Z^*Z)=|\det Z|^2$.  The quaternionic positivity criterion is
\[
 G\succeq0,
 \quad\Ran Z\subset\Ran G,
 \quad\lambda\ge\tr(Z^*G^\dagger Z).
\]
The complex block criterion is
\[
 G\succeq0,
 \quad\Ran Z\subset\Ran G,
 \quad\lambda I_2-Z^*G^\dagger Z\succeq0.
\]
Since a positive $2\times2$ matrix has largest eigenvalue at most its trace, the first scalar inequality implies the second matrix inequality.  The example has complex Schur complement zero but violates the scalar trace bound.
\end{proof}

\begin{remark}[Branch count]
On the diagonal locus $Z=0$ and $G=\diag(x,y)$, the quaternionic cubic is
$\lambda xy$, whereas the complex quartic is $\lambda^2xy$.  Thus the thirteen-dimensional quartic model merges two of four diagonal branches into one multiplicity-two branch.  Four independently variable simple diagonal branches require a second scalar coordinate, giving at least the fourteen-dimensional block
$\diag(x,y,\lambda,\mu)$.  This is the exact branch morphism behind the distinction; the equation $z_0z_1z_2z_3=t$ is not produced merely by writing $4+4+4+1$.
\end{remark}

\section{Exact relation to the existing colour \texorpdfstring{$S^{12}$}{S12}}

Let
\begin{equation}\label{eq:colorV}
 V_{\mathrm{col}}=\Mat_{3\times2}(\C),
 \qquad
 \Gamma(A)=A^*A\in\Herm_2(\C)_{\succeq0}.
\end{equation}
The left $SU(3)$ Gram quotient is
$V_{\mathrm{col}}/SU(3)\cong\Herm_2(\C)_{\succeq0}$.
Define the proper map
\begin{equation}\label{eq:color-to-13}
 F_{13}(A)
 =X\begin{pmatrix}\Gamma(A)&0\\0&0\end{pmatrix}X^*
 \in\partial\cC_{13,P}.
\end{equation}

\begin{proposition}[The two $S^{12}$ carriers are connected by a typed map]\label{prop:S12-map}
The map $F_{13}$ factors exactly as
\begin{equation}\label{eq:S12-factor}
 \Mat_{3\times2}(\C)
 \xrightarrow{\ \Gamma\ }
 \Herm_2(\C)_{\succeq0}
 \xrightarrow{\ \jmath_P\ }
 \partial\cC_{13,P}.
\end{equation}
After one-point compactification it gives
\begin{equation}\label{eq:S12-compact-map}
 \boxed{
 S^{12}_{\mathrm{col}}
 \xrightarrow{\ \Gamma^+\ }
 B^4
 \xrightarrow{\ \jmath_P^+\ }
 S^{12}_{13}.}
\end{equation}
This sphere map is based null-homotopic because it factors through the based-contractible ball $B^4$.  Its image is the embedded four-ball face $\jmath_P^+(B^4)$; it is not an identification of the two sphere constructions.
\end{proposition}

\begin{proof}
The factorization is the definition.  Properness follows from
$\tr\Gamma(A)=\|A\|_F^2$ and injectivity/properness of $\jmath_P$ on the cone.  The compactified Gram quotient has codomain $B^4$, and a closed ball contracts to its chosen boundary basepoint while fixing that point.
\end{proof}

\begin{remark}[Relation to the toroidal map]
The toroidal cusp construction already uses the same marked integral matrix $P$ as the angular part of the lattice map $\Phi=P\oplus Q$.  The new quaternionic coordinate $w$ has not been assigned an integral period lattice, monodromy, fan ray, or filling homomorphism.  Therefore the exact new maps are \eqref{eq:pullback-square}, \eqref{eq:iota-cones}, and \eqref{eq:S12-factor}; a map from the completed toroidal total space to $S^{12}_{13}$ requires additional transition and filling data and is not inferred from the shared symbol $P$.
\end{remark}

\section{A precise threshold consequence of the cubic}

Let $G\in\Herm_2(\C)_{>0}$ and retain a conversion constant $\alpha>0$.  Define
\begin{equation}\label{eq:momentum}
 p=\alpha G,
 \qquad m^2=\det p=\alpha^2\det G.
\end{equation}
In Dirac variables put
\begin{equation}\label{eq:QD}
 Q_G(\psi_L,\psi_R)
 =\psi_L^*\adj(G)\psi_L+\psi_R^*G\psi_R\ge0.
\end{equation}
Then
\begin{equation}\label{eq:N-threshold}
 \cN_{13}=\lambda\det G-Q_G.
\end{equation}

\begin{proposition}[Joint scalar--cubic threshold]\label{prop:threshold}
Fix $\lambda_0>0$ and $\nu_0>0$.  On the invariant level set
\begin{equation}\label{eq:level-set}
 \lambda=\lambda_0,
 \qquad \cN_{13}=\nu_0,
 \qquad G>0,
\end{equation}
one has the exact lower bound
\begin{equation}\label{eq:mass-threshold}
 \boxed{m^2\ge\frac{\alpha^2\nu_0}{\lambda_0}.}
\end{equation}
Equality holds exactly when $\psi_L=\psi_R=0$.  At a fixed $G$ above threshold, the spinor fibre is the ellipsoid
\begin{equation}\label{eq:S7-fibre}
 Q_G(\psi_L,\psi_R)
 =\lambda_0\det G-\nu_0,
\end{equation}
which is diffeomorphic to $S^7$; it collapses to a point at threshold.
\end{proposition}

\begin{proof}
Equation \eqref{eq:N-threshold} gives
\[
 \lambda_0\det G=\nu_0+Q_G\ge\nu_0.
\]
Since both $G$ and $\adj(G)$ are positive definite, equality holds precisely at the zero Dirac spinor.  Above threshold, $Q_G$ is a positive-definite quadratic form on $\C^4\cong\R^8$, and a positive level is an ellipsoidal $S^7$.
\end{proof}

\begin{remark}
This is a finite-dimensional invariant threshold, not a quantum-field-theoretic mass gap.  A Hamiltonian statement would additionally require one fixed positive Hilbert space, a self-adjoint dynamics, a vacuum projection, and an energy-controlled map from these variables into physical states.
\end{remark}

\section{Conclusion}

The expression $4+4+4+1$ has a rigid continuation of the preceding marked construction.  Extend the same observation map $P$ from complex to quaternionic Hermitian matrices, but require its observed block to remain in the original four-dimensional Lorentz space $\Herm_2(\C)$.  The result is the thirteen-dimensional pullback $\cJ_{13,P}$, with one Lorentz vector, two Weyl spinors (equivalently one Dirac spinor), and one scalar.  Its cubic restricts exactly to the nine-dimensional cubic, its positive cone has compactified rank defect $S^{12}$, and the old $S^8$ is the fixed unknotted subspace with quaternionic linking sphere $S^3$.

The same coordinate count also supports a complex rank-four determinant, but the exact identity
\[
 \det\cK_{13}=\lambda\cN_{13}+|\det Z|^2
\]
separates the two constructions.  The quaternionic slice preserves the cubic ladder; the complex slice carries the quartic and Pluecker correction.  Finally, Jordan multiplication generates two further Lorentz-vector coordinates and closes the thirteen-dimensional slice to the full fifteen-dimensional algebra $\Herm_3(\Hh)$.  These maps specify what the thirteen-dimensional object is, how it contains the old object, how it connects to the existing colour $S^{12}$, and exactly where further toroidal or physical data would have to enter.

\appendix
\section{Symbolic replay}

The accompanying exact-arithmetic script verifies:
\begin{itemize}
\item all retained $P,R_P,u,X$ identities;
\item the quaternionic splitting $\psi^*A\psi=z^*Az+w^*Aw$;
\item $\det_{\C}\rho(H)=\cN_{13}(H)^2$;
\item the Dirac charge-conjugation identity;
\item the adjugate covariance and internal $SU(2)$ invariance;
\item exact restriction to the old cubic and the determinant correction under $r_9$;
\item the complex rank-four determinant and the Pluecker correction;
\item the explicit Jordan nonclosure witness.
\end{itemize}

\end{document}
